\subsection{Central Frobenius over a general CM center}

The imaginary quadratic calculation extends to CM centers by adding
a unit-norm obstruction. The relevant ideal and unit groups are the
ones used in polarized ideal-class theory; compare
\citet[Lemma 5.5 and \S5.3]{BrooksEtAl2017}. We give the argument
directly for the descent question.

\begin{proposition}\label{prop:central-cm}
Let $A/\overline{\F}_p$ be a geometrically simple abelian variety, with geometric
endomorphism division algebra $D$ of degree $d$ over a CM field
$K$. Choose a model over $\F_{p^N}$ on which all geometric
endomorphisms are defined, with Frobenius $\pi_0\in K$. For
$v\mid p$, write $e_v=\ord_v(p)$, let $f_v$ be the residue degree,
and put $\lambda_v=\ord_v(\pi_0)/(Ne_v)$. Define
\[
 a=\operatorname{lcm}_{v\mid p}\operatorname{den}(e_v\lambda_v),
 \qquad
 I=\prod_{v\mid p}\pp_v^{a e_v\lambda_v},
 \qquad t=\ord_{\Cl(K)}[I],
\]
where $\operatorname{den}$ is the positive reduced denominator.
Choose $(\eta)=I^t$, and let
\[
 u=\frac{\eta\bar\eta}{p^{at}}\in\OO_{K^+}^{\times,+},
 \qquad
 s=\ord\left([u]\in
 \OO_{K^+}^{\times,+}/N_{K/K^+}(\OO_K^\times)\right).
\]
Here $K^+$ is the maximal totally real subfield and the superscript
$+$ means totally positive. Then $s\in\{1,2\}$, independently
of the choice of $\eta$. The possible central-Frobenius model
exponents are exactly the positive multiples of
\begin{equation}\label{eq:cm-central-min}
                         n_{\mathrm{cent}}=ats.
\end{equation}
Moreover
\[
 a\mid d\operatorname{lcm}_{v\mid p}f_v,
 \qquad n_{\mathrm{cent}}\mid 2h_Ka
               \mid 2h_Kd\operatorname{lcm}_{v\mid p}f_v.
\]
If $D=K$, formula \eqref{eq:cm-central-min} is the minimum among
all models in the geometric isogeny class.
\end{proposition}

\begin{proof}
The Weil relation gives $\lambda_v+\lambda_{\bar v}=1$ and
$e_v=e_{\bar v}$, whence $I\bar I=(p)^a$. Thus the displayed
$u$ is a totally positive unit of $K^+$. Replacing $\eta$ by
$\eta\varepsilon$ multiplies $u$ by the norm of a unit. For
every totally positive unit $z$ of $K^+$ one has
$z^2=N_{K/K^+}(z)$. Consequently the unit quotient has exponent
at most two, proving the assertions about $s$.

Let $\beta\in K$ be the Frobenius of a compatible central model
over $\F_{p^n}$, transported into $D$ by a geometric isogeny.
Over a common finite field, powers of $\beta$ and $\pi_0$ agree.
Their slopes are therefore equal under this identification of
$K$. The integer valuations $\ord_v(\beta)=ne_v\lambda_v$
force $a\mid n$. Write $n=am$. Then $(\beta)=I^m$, so
$t\mid m$. Writing $m=tr$, there is a unit $\varepsilon$ with
$\beta=\eta^r\varepsilon$. The Weil condition is precisely
\[
       \beta\bar\beta=p^{atr}
       \quad\Longleftrightarrow\quad
       u^rN_{K/K^+}(\varepsilon)=1.
\]
It is possible exactly when $s\mid r$. This proves necessity.

Conversely choose a unit $\varepsilon$ with
$N_{K/K^+}(\varepsilon)=u^{-s}$ and set
$\beta=\eta^s\varepsilon$, $n=ats$. The ideal of $\beta$ is
integral and $\beta\bar\beta=p^n$, so $\beta$ is a Weil
$p^n$-number. The quotient $\beta^N/\pi_0^n$ has valuation zero
at every finite prime and absolute value one at every complex
embedding. It is a root of unity, say of order $b$. Hence
$\beta^{Nb}=\pi_0^{nb}$.

Every positive power of $\pi_0$ generates $K$: all endomorphisms
remain defined after extension, so Honda--Tate identifies the
center there with $\Q(\pi_0^r)=K$. The equality of powers now
also gives $\Q(\beta^r)=K$ for every $r>0$. The slopes of
$\beta$ have the same local invariants as those of $\pi_0$.
Honda--Tate therefore gives the same division-algebra degree and
dimension, and the equality of powers gives a geometric isogeny.
This constructs the central model at $n=ats$; base extension gives
all its positive multiples.

Finally $[K_v:\Q_p]=e_vf_v$, and the local invariant of $D$ is
$e_vf_v\lambda_v$ modulo $\Z$. Its denominator divides $d$.
Thus the denominator of $e_v\lambda_v$ divides $df_v$, proving
the divisibility for $a$. Now $t\mid h_K$ and $s\mid2$ give
the two remaining divisibilities. If $D=K$, every transported
Frobenius is central, so there are no additional models to consider.
\end{proof}

For an imaginary quadratic center with split $p$ as in
Proposition~\ref{prop:central}, one has $a=d$,
$I=\pp^j\bar\pp^{d-j}$, and $K^+=\Q$. The positive unit
quotient is trivial, while
\[
 \ord[I]=\frac{\ord[\pp]}{\gcd(\ord[\pp],2j-d)}.
\]
Thus the general formula recovers \eqref{eq:central-min}.
For a higher-degree center, the residue degrees and the unit-norm
class are explicit additional data. The bound just proved is not
asserted to be sharp. When $D$ is noncommutative, a noncentral
Frobenius may give a smaller exponent; the formula does not exclude
that possibility.
